\section{Síntesis y ampliación}

Esta clase devuelve el RL del terreno de los algoritmos de optimización al de los sistemas de producto. La product belief determina qué conducta se quiere ver; la interface expone las tareas y los fallos; la eval descompone la «sensación» en evidence reproducible; el data debugging rastrea el origen de las conductas; el RL environment define estados, acciones, herramientas y recuperación; la reward incorpora en el objetivo tanto los resultados como el proceso; y la evidence de los usuarios tras la puesta en producción inicia la siguiente ronda de investigación. Si falta cualquiera de estas capas, otra capa tendrá que cargar con problemas que no puede resolver.

\subsection{Una tabla que articula el co-design loop}

\begin{table}[H]
\centering
\small
\begin{tabular}{@{}L{0.16\textwidth}L{0.24\textwidth}L{0.27\textwidth}L{0.23\textwidth}@{}}
\toprule
Capa & Objeto central & Fallo representativo & Pregunta de aceptación \\
\midrule
Belief & Conducta objetivo y relación con el usuario & Lemas que no pueden llevarse a la práctica & ¿Puede escribirse como un contract de conducta? \\
Interface & artifact, diff, control y permisos & Capacidad invisible o irreversible & ¿El usuario puede entender, editar y detener? \\
Eval & Distribución, rubric, judge, umbrales & Desajuste entre el benchmark y las tareas reales & ¿Es reproducible, está dividida por slice y permite decidir? \\
Data & provenance, mixture, preference & Conflictos y contaminación de datos & ¿Puede localizarse el origen de una conducta y revertirse? \\
Environment & Estados, acciones, herramientas, terminación & Las tareas de entrenamiento difieren de las de despliegue & ¿Coinciden permisos, presupuesto y recuperación ante fallos? \\
Reward & outcome, process, cost, risk & reward hacking / Goodhart & ¿El verifier es independiente y resistente a ataques? \\
Deployment & telemetry, regresiones, responsables & El éxito de la demo oculta incidentes de cola larga & ¿Se monitorea el outcome y puede revertirse con rapidez? \\
\bottomrule
\end{tabular}
\caption{Mapa de aceptación en siete capas para el diseño conjunto de RL, producto e investigación.}
\end{table}

\subsection{Preguntas que deben responderse antes de iniciar un proyecto de Frontier Product}

\begin{importantbox}{Autoevaluación en doce preguntas}
\begin{enumerate}
  \item ¿La product belief puede escribirse como conducta observable y no como adjetivos?
  \item ¿Cuáles son las tareas reales del usuario objetivo y el costo de los fallos?
  \item ¿Mediante qué form factor se expone la capability del modelo?
  \item ¿Qué artifacts, estados y fuentes deben poder rastrearse?
  \item ¿Qué facultades de confirmar, editar, detener y hacer rollback conserva el usuario?
  \item ¿Están fijos la distribución, el judge, la rubric y la versión de la eval?
  \item ¿Se reportan a la vez los resultados overall, por slice, worst-case y los casos graves?
  \item ¿La taxonomy de fallos puede asociarse con causas raíz en los datos, el sistema o la policy?
  \item ¿Las herramientas, los permisos, el presupuesto y la terminación del RL environment coinciden con los del despliegue?
  \item ¿La reward separa el outcome real, el proceso, el costo y el riesgo?
  \item ¿Qué blind spots tiene el verifier y puede la política atacarlo?
  \item Tras la puesta en producción, ¿quién revisa el outcome, quién puede revertir y qué condiciones activan la detención del servicio?
\end{enumerate}
\end{importantbox}

\subsection{Conjunto mínimo de registros para la reproducibilidad}

Ya se investiguen data, behavior, RL environment o interface, deben guardarse source/version, prompt/system, el modelo y el sampling, el schema de las herramientas, la trace completa, judge/rubric, el disagreement humano, las muestras de fallo y las condiciones de reversión. Si se usa un AI judge o synthetic data, también deben fijarse el modelo generador, el modelo judge y los permisos; de lo contrario, un resultado no podrá ser explicado por la versión siguiente.

\begin{lstlisting}[caption={Registro reproducible de un experimento de Frontier product}]
experiment = {
    "hypothesis": product_belief_as_testable_behavior,
    "model": [base, post_training, checkpoint, sampling],
    "environment": [tools, permissions, budget, termination],
    "interface": [surface, artifact_schema, confirmation, rollback],
    "evaluation": [datasets, slices, judges, rubrics, thresholds],
    "evidence": [traces, outcomes, disagreements, severe_failures],
    "decision": [ship, hold, rollback, next_hypothesis],
}
\end{lstlisting}

\subsection{Lecturas adicionales}

\begin{itemize}
  \item Bai et al., \emph{Constitutional AI: Harmlessness from AI Feedback}: para entender RLAIF, las restricciones basadas en principios y la AI feedback.
  \item Röttger et al., \emph{XSTest}: estudia la exaggerated safety y el over-refusal ante benign prompts.
  \item Lilian Weng, \emph{Reward Hacking in Reinforcement Learning}: recopila de forma sistemática la proxy reward y las vías de hacking.
  \item Ouyang et al., \emph{Training language models to follow instructions with human feedback}: el pipeline clásico de RLHF y sus limitaciones.
  \item Amodei et al., \emph{Concrete Problems in AI Safety}: para entender el diseño de entornos a partir del reward hacking, los side effects y la oversight.
\end{itemize}

\begin{warningbox}{Límite final}
Los nombres de productos, el desempeño de los modelos, las curvas de costo y las prácticas organizacionales de esta clase corresponden al momento en que se impartió, en 2025. Las notas conservan el método de co-design; no tratan ninguna demo, benchmark o proceso de empresa en particular como un hecho permanente. La capacidad realmente transferible consiste en traducir la visión en un entorno, la conducta en una eval, rastrear los fallos hasta los datos y el sistema, y luego decidir con evidencia reproducible si se entrega.
\end{warningbox}

\end{document}
